\section{Failure-First Action-Pool Audit}
\label{sec:audit}

\subsection{Composite Actions and Strict Utility}

The single-best/virtual-best gap is established in per-instance algorithm
selection \cite{Rice1976AlgorithmSelection,LindauerEtAl2019AlgorithmSelection}
and has a direct LLM-routing analogue in LLMRouterBench's Oracle analysis
\cite{LiEtAl2026LLMRouterBench}. We operationalize that comparison for composite
RAG actions, then place it inside a broader audit that also checks interfaces,
absolute capability, physical-cost comparability, and validation feasibility.

Let $\mathcal{A}$ be the actions available to a router and let $q_i$ denote one of
$N$ evaluation queries. An action fixes the generator, retrieval mechanism, memory
organization, prompt contract, and deterministic guard. Its strict outcome is
$S_{ia}\in\{0,1\}$. Because the outcome includes answer, evidence, abstention, and
interface requirements, two routes that produce the same answer string need not
have the same utility.

The best fixed utility of an action pool is

\begin{equation}
B(\mathcal{A})=\max_{a\in\mathcal{A}}
\frac{1}{N}\sum_{i=1}^{N}S_{ia}.
\label{eq:best-fixed}
\end{equation}

A label-aware oracle that chooses any successful action separately for each query
has utility

\begin{equation}
O(\mathcal{A})=\frac{1}{N}\sum_{i=1}^{N}
\max_{a\in\mathcal{A}}S_{ia}.
\label{eq:oracle}
\end{equation}

We refer to the difference as \routingheadroom{}:

\begin{equation}
H(\mathcal{A})=O(\mathcal{A})-B(\mathcal{A}).
\label{eq:headroom}
\end{equation}

$O$ is an upper bound derived from outcomes, not a policy that can be executed
without labels. If $O$ is below a required utility $q_{\min}$, the action pool is
\emph{capability-infeasible} for that requirement. If $H=0$, no selector can improve
strict success over the best fixed action on the observed matrix. Positive headroom
is necessary but not sufficient: a practical policy must still predict the useful
differences from information available before generation and generalize beyond the
matrix.

\subsection{Cost Structure}

For route $a$, let $E_{ia}$ be the selected-board generation energy when telemetry
is valid. Mean cost is $\bar E_a=N^{-1}\sum_iE_{ia}$. An action $a$ descriptively
dominates $b$ in the observed quality--energy plane when
$\bar S_a\geq\bar S_b$ and $\bar E_a\leq\bar E_b$, with at least one strict
inequality. Physical interpretation additionally requires a common measurement
boundary and attention to device assignment. In this study, all traces use the
same generation-window definition, but not every pair ran on the same T4 board;
cross-board dominance is therefore descriptive rather than causal.

An energy-aware router can have value even when its success headroom is zero if
multiple actions cover the same successful queries at meaningfully different
costs. If the best fixed action covers every successful query and also has the
lowest route-level mean energy, then no other fixed route improves the observed
mean-quality/mean-energy point. This statement does not exclude small per-query
cost substitutions caused by trace-level energy variation; we audit that
label-aware bound separately rather than equating route diversity with a useful
Pareto frontier.

\subsection{Five Audit Questions}

The proposed audit proceeds from route validity to policy feasibility:

\begin{enumerate}
    \item \textbf{Interface compatibility:} do retained model--prompt--runtime
    combinations terminate and produce parseable, support-checkable outputs?
    \item \textbf{Capability:} does $O(\mathcal{A})$ meet the declared operating
    requirement, or is the action set inadequate even with label access?
    \item \textbf{Complementarity:} is $H(\mathcal{A})>0$, and which actions add
    unique successful queries?
    \item \textbf{Cost comparability:} are costs observed under an explicit common
    boundary, and do successful actions expose a genuine quality--cost choice?
    \item \textbf{Constrained selectability:} can a decision-time policy satisfy
    validation requirements for success, coverage, latency, and abstention?
\end{enumerate}

Failure at an early question changes the interpretation of later stages. An
incompatible action should be repaired under a new specification or removed before
expensive matrix collection. A capability-infeasible pool should be expanded before
router fitting. Zero complementarity distinguishes action-space failure from
classification error. A measurement qualification limits energy claims. Finally,
validation infeasibility should produce an explicit non-deployment state rather
than a threshold chosen after inspecting test labels.

Figure~\ref{fig:routability-audit} summarizes these questions and the observed
status of the resident-eligible pool. Its statuses are retrospective diagnostics,
not prospectively frozen gates for this experiment.

\begin{figure*}[t]
    \centering
    \includegraphics[width=0.96\textwidth]{routability_audit.png}
    \caption{The five-question audit applied retrospectively to the frozen
    resident-eligible pool. The
    statuses summarize existing traces: three eligible Granite-1B actions never
    parsed; even a per-query oracle reached only 12.5\% against the declared 80\%
    requirement; routing headroom was zero; energy telemetry was complete within
    the declared generation window but spanned physical boards; and no frozen
    threshold was validation-feasible.}
    \label{fig:routability-audit}
\end{figure*}

\subsection{Confirmatory and Diagnostic Roles}

The learned controller and its hypothesis test were frozen before test-label
access. The oracle ceilings, headroom, interaction contrasts, task-aware oracle,
and failure decompositions were computed afterward from the complete sealed matrix
in response to the negative result. We label them post-hoc diagnostics throughout.
They explain why the confirmatory policy failed; they do not convert that failure
into a successful prospective router claim.
